\chapter{Introduction}
\label{ch:intro}

\section{Context of the Research}
Integrating artificial intelligence into medical imaging, particularly through deep learning and convolutional neural networks (CNNs), represents a profound paradigm shift in modern dentistry \cite{alenezi2023artificial, bozkurt2023deep, khan2021clinical, lecun2015deep}. Over the past decade, the rapid advancement of computer vision architectures—most notably the "You Only Look Once" (YOLO) family of object detection models—has demonstrated exceptional capability in analyzing complex anatomical structures in real-time \cite{geetha2024ai, lin2014microsoft, redmon2016you}. Dental radiography, including panoramic and bitewing X-rays, provides critical diagnostic information; however, interpreting these multimodal images is often subjective and dependent on the practitioner's clinical experience \cite{chen2022comprehensive, kim2020automated, rahman2021deep}. Recent studies have highlighted the efficacy of advanced object detection algorithms, such as YOLOv8, in autonomously identifying and localizing various dental anomalies, including interproximal caries, periapical lesions, and impacted teeth \cite{devito2023advanced, hao2023innovative, lee2023detection, sharma2024early}. By bridging the gap between raw radiological data and actionable diagnostic insights, AI-assisted computer vision serves as an invaluable second-opinion tool in contemporary endodontic and restorative workflows \cite{endres2020development, hwang2019overview, mohammad2022deep, wang2022automated}.

\section{Problem Definition}
Despite the promising capabilities of deep learning in medical diagnostics, standardizing automated dental anomaly detection remains a highly complex scientific challenge \cite{haghanifar2020dental, jocher2023ultralytics, nguyen2023dental, schwendicke2020artificial}. The primary difficulty lies in the high variability of dental X-rays, where issues such as overlapping teeth structures, low image contrast, artifacts from metallic restorations, and diverse anatomical morphologies obscure subtle pathologies like early-stage caries or discrete periapical lesions \cite{lee2021challenges, orhan2020evaluation, tuzoff2019tooth}. Furthermore, existing AI models frequently struggle with class imbalance and inconsistent bounding box precise localization when trained on heterogeneous datasets from varying clinical environments \cite{pang2019libra, prados2020artificial, son2022performance}. While older generalized neural networks fall short of offering real-time, multi-class diagnostic precision without exhausting computational resources, state-of-the-art single-stage detectors need severe custom architectural optimization and properly mapped class ontologies to correctly differentiate between visually similar dental conditions \cite{tan2019efficientnet, xie2023survey}. Consequently, there is a critical need for an adapted, highly precise, and computationally efficient deep learning pipeline that can reliably operate on local clinical hardware without sacrificing diagnostic accuracy.

\section{Proposed Solution and Motivation}
This work investigates the development and implementation of an optimized, AI-assisted dental imaging web platform driven by a customized YOLOv8 object detection architecture. The primary objective is to train a robust neural network on a unified, comprehensively annotated dataset of dental radiographs to autonomously detect, classify, and visualize pathological findings with high diagnostic confidence. By integrating this advanced computer vision model into a secure and intuitive web interface, the research aims to provide a seamless, real-time diagnostic support tool for clinical environments. The motivation behind this endeavor is deeply rooted in enhancing diagnostic precision and reducing physician fatigue; by providing dentists with an instantaneous visualization of potential anomalies categorized by standardized classes, the proposed system aspires to significantly minimize human error, expedite treatment planning, and ultimately improve the overarching quality of patient care in modern dentistry.